% SemanticLongTail — technical report (skeleton)
% Open source at https://github.com/ceurws/SemanticLongTail
%
% Target: technical report on Zenodo. Not peer-reviewed.
% Discussion: https://github.com/ceurws/semantify/discussions/76
%
% CEUR-WS 1-column style (ceurart). Do NOT use the twocolumn / hf options.
% The report follows the CEUR-WS Policy on AI-Assisting Tools although it is
% not published via CEUR-WS; see AGENTS.md.
%
% Sections are intentionally EMPTY until the authors write them.
% The report must be self-contained: every term, concept, walk and method
% step is defined in the text; links only point to further material.
%
\documentclass{ceurart}
% typographical quoting
\newcommand{\tquote}[1]{``#1''}
% footnote links
\newcommand{\footnotelink}[2]{\href{#1}{#2}\footnote{#1}}
\sloppy

\begin{document}

%% Rights management information (CC-BY is the CEUR-WS default)
\copyrightyear{2026}
\copyrightclause{Copyright for this paper by its authors.
  Use permitted under Creative Commons License Attribution 4.0
  International (CC BY 4.0).}

%% Conference information
\conference{Technical report, 2026}

%% Title
% Working title, to be finalized by the authors.
\title{The Semantic Long Tail: what 25 years of CEUR-WS volume pages say beyond their tags}
\tnotemark[1]
\tnotetext[1]{Technical report. This version has not been peer-reviewed.}

%% Authors
% Confirmed authors only. ORCIDs taken from the media wiki scholar pages.
% Further co-authors are NOT listed until participation is confirmed.
% AI agents are never authors (CEUR-WS HOWTOSUBMIT FAQ 15).
\author[1,2]{Wolfgang Fahl}[%
orcid=0000-0002-0821-6995,
email=wf@bitplan.com,
]
\cormark[1]
\address[1]{RWTH Aachen University, Computer Science i5, Aachen, Germany}
\address[2]{BITPlan GmbH, Willich, Germany}
\cortext[1]{Corresponding author.}

%% Abstract (author-provided text).
%% The abstract environment of ceurart is verbatim: its first line must be
%% text, a comment-only first line makes it miss \end{abstract}.
\begin{abstract}
Abstract to be written by the authors.
\end{abstract}

%% Keywords
\begin{keywords}
% TODO(authors)
\end{keywords}

\maketitle

\section{Introduction}\label{sec:introduction}
% TODO(authors)

\section{Background}\label{sec:background}
% TODO(authors): the CEUR-WS volume page and its tags, the Vol-XXX template
% and its versions, the three layers marked, conventional and deduced,
% graph walk and PentaLink; define every term used later.

\section{Prior Art}\label{sec:priorart}
% TODO(authors): the semantification papers of 2023 and 2026, the SWJ
% submission of 2024 and its reviews, prior art on graph walks and
% PentaLinks, related conversions of proceedings pages.

\section{Research Question}\label{sec:rq}
% TODO(authors): the open question in the authors' own words; why the
% concepts are found, not given.

\section{Data}\label{sec:data}
% TODO(authors): the population of volume pages and papers, the snapshot
% date, the share of volumes per template version.
% No CEUR-WS pages or PDFs are redistributed; see data/README.md.

\section{Method}\label{sec:method}
% TODO(authors): the residue, clustering and counting; the graph walk per
% concept and the forks it shows; competency questions as named queries
% against the endpoint; sharpening of the ontology per round; the agentic
% work and the provenance notation of the repository are described here,
% in the text.

\section{Concepts Found}\label{sec:concepts}
% TODO(authors): one subsection per round, each with the concept, its walk,
% its counts, the questions it raised, the ontology change and what the
% round left unexplained.

\section{Disambiguation Needs}\label{sec:disambiguation}
% TODO(authors): the forks counted over all volumes: slug with several
% pids, slug with no pid, pid depending on time.

\section{Production State}\label{sec:production}
% TODO(authors): the state of the tooling on ceur-ws.org at a stated date;
% the frozen release the report refers to.

\section{Threats to Validity}\label{sec:threats}
% TODO(authors)

\section{Conclusion}\label{sec:conclusion}
% TODO(authors)

%% Declaration on Generative AI (required by CEUR-WS, see
%% https://ceur-ws.org/GenAI/Policy.html and https://ceur-ws.org/GenAI/Taxonomy.html).
%% Finalized by the authors from PROVENANCE.md before release.
%% Must name every tool and map each contribution to a taxonomy category.
%% Agents used as research instruments for the method are described in
%% Section Method and named here as such, not as writing assistance.
\section*{Declaration on Generative AI}
% TODO(authors)

%% Bibliography
\nocite{ceurws2024genaipolicy,ceurws2024genaitaxonomy}
\bibliography{references}

\end{document}
